% Condensed "What Is Biomedical Text?" (EN). One figure: the discharge summary
% (signature). The article/forum registers are evoked in prose (15-page limit).

\noindent
\begin{center}
\begin{tikzpicture}
\node[draw=ThesisNeutral, fill=ThesisPaper, rounded corners=3pt, line width=0.5pt,
      inner sep=10pt] {%
\begin{tabular}{@{}l !{\hspace{6pt}\textcolor{ThesisNeutral}{\vrule width 0.3pt}\hspace{6pt}} l@{}}
\begin{minipage}[t]{0.40\textwidth}
\scriptsize\color{ThesisInk}
{\tiny\sffamily\color{ThesisNeutral} FRANÇAIS}\\[4pt]
{\bfseries CENTRE HOSPITALIER DE SAINT-AUBIN}\\ Dossier n\textsuperscript{o} 4471902\\[1pt]
Cardiologie, Soins Intensifs\\ Séjour : 03/11--08/11/2024\\[4pt]
{\color{ThesisNeutral}\rule{\linewidth}{0.3pt}}\\[4pt]
\textbf{Patient :} CARON Julien, né le 14/03/1962 (62 ans). \textbf{Dest. :} D\textsuperscript{r} S.~Lemaire.\\[5pt]
\textbf{Motif.} Douleur thoracique constrictive rétrosternale, irradiation bras G, depuis 2h.\\[5pt]
\textbf{Antécédents.} HTA. DT2 sous metformine. Tabac 30 PA. Dyslipidémie.\\[5pt]
\textbf{Examen.} TA 148/92, FC 98, SpO\textsubscript{2} 96\% AA, EVA 7/10. Pas de signe d'IVG.\\[5pt]
\textbf{Paraclinique.} ECG : sus-décalage ST V1--V4. Tropo T 4520 ng/L (N $<$14). Coro : occlusion IVA proximale, angioplastie + stent actif, TIMI 3. ETT : FEVG 40\%.\\[5pt]
\textbf{Au total.} SCA ST+ antérieur (I21.0).\\[5pt]
\textbf{Ttt sortie.} Kardégic 75, ticagrélor 90x2, atorvastatine 80, bisoprolol 2,5, ramipril 5.\\[5pt]
\textbf{CAT.} Réadaptation cardiaque. Consultation cardio à 1 mois.\hfill D\textsuperscript{r} C.~Després
\end{minipage}
&
\begin{minipage}[t]{0.40\textwidth}
\scriptsize\color{ThesisNeutral}
{\tiny\sffamily\color{ThesisNeutral} ENGLISH · TRANSLATION}\\[4pt]
{\bfseries SAINT-AUBIN HOSPITAL}\\ Record no. 4471902\\[1pt]
Cardiology, Intensive Care\\ Stay: 03/11--08/11/2024\\[4pt]
{\color{ThesisNeutral}\rule{\linewidth}{0.3pt}}\\[4pt]
\textbf{Patient:} CARON Julien, b.\ 14/03/1962 (62 y). \textbf{To:} Dr S.~Lemaire.\\[5pt]
\textbf{Reason.} Constrictive retrosternal chest pain, radiating to L arm, for 2h.\\[5pt]
\textbf{History.} HTN. T2DM on metformin. Smoking 30 PY. Dyslipidaemia.\\[5pt]
\textbf{Exam.} BP 148/92, HR 98, SpO\textsubscript{2} 96\% RA, pain 7/10. No sign of LVF.\\[5pt]
\textbf{Workup.} ECG: ST elevation V1--V4. Trop T 4520 ng/L (N $<$14). Angiography: proximal LAD occlusion, angioplasty + drug-eluting stent, TIMI 3. Echo: LVEF 40\%.\\[5pt]
\textbf{Summary.} Anterior STEMI (I21.0).\\[5pt]
\textbf{Discharge tx.} Aspirin 75, ticagrelor 90x2, atorvastatin 80, bisoprolol 2.5, ramipril 5.\\[5pt]
\textbf{Plan.} Cardiac rehabilitation. Cardiology review at 1 month.\hfill Dr C.~Després
\end{minipage}
\end{tabular}%
};
\end{tikzpicture}
\captionof{figure}{A discharge summary for an acute myocardial infarction; French original, English translation.}
\label{fig:bt-clinical-note}
\end{center}
\medskip

\noindent \Cref{fig:bt-clinical-note} is a letter from one cardiologist to a colleague. For an
ordinary French reader it is very hard to read: few verbs, telegraphic sentences,
information dense and simply separated by commas like a table. This is not
carelessness but a style of communication specially tuned for the community it is
written for. The linguist Zellig Harris called such a specialised, rule-governed
form a \emph{sublanguage}~\citep{harris_mathematical_1968}, and the clinical note
is its best-documented case~\citep{sager_nlip_1981}. It is above all a form of
compression: \citet{anderson_compression_1975} showed that clinical shorthand
obeys a small set of deletion rules, dropping only what the reader can put back.
\emph{Chest films normal} stands for \emph{the chest films were normal}; what
disappears is the part the expert reader already knows~\citep{stalnaker_common_2002}.
The note is also performative: a prescription lets the patient obtain the
medication, an order to transfer a patient sets the transfer in
motion~\citep{austin_words_1962}.

All of this makes the note worth studying, but it also makes it hard to obtain. The
note rarely leaves the hospital that produced it, because it is full of personal
data, and removing the obvious identifiers does not fix this: a pseudonymised note
is still not anonymous. For French the problem is worse. There is no shareable
corpus of real French hospital records~\citep{neveol_clinical_2018}. Public data is
an alternative: published case reports, drug leaflets, article abstracts. But a
case report is written to teach; it reads like a short story, in full sentences and
the past tense. The same infarction, told in a research article or on a patient
forum (\Cref{fig:bt-web}), is abundant and public, yet it speaks a different language.

\begin{center}
\begin{tikzpicture}
\node[draw=ThesisNeutral, fill=ThesisPaper, rounded corners=3pt, line width=0.5pt,
      inner sep=10pt] {%
\begin{tabular}{@{}l !{\hspace{6pt}\textcolor{ThesisNeutral}{\vrule width 0.3pt}\hspace{6pt}} l@{}}
\begin{minipage}[t]{0.40\textwidth}
\footnotesize\color{ThesisInk}
{\tiny\sffamily\color{ThesisNeutral} FRANÇAIS}\\[5pt]
\textit{Bonjour, mon père a fait une crise cardiaque la semaine dernière. Il a eu
une grosse douleur dans la poitrine qui descendait dans le bras gauche, il était
tout pâle et en sueur, on a appelé le 15 tout de suite. Est-ce que quelqu'un sait
combien de temps dure la convalescence après un infarctus ? Merci d'avance.}
\end{minipage}
&
\begin{minipage}[t]{0.40\textwidth}
\footnotesize\color{ThesisNeutral}
{\tiny\sffamily\color{ThesisNeutral} ENGLISH · TRANSLATION}\\[5pt]
\textit{Hello, my father had a heart attack last week. He had a strong pain in the
chest that went down into the left arm, he was all pale and sweating, we called
emergency services right away. Does anyone know how long recovery takes after a
heart attack? Thanks in advance.}
\end{minipage}
\end{tabular}%
};
\end{tikzpicture}
\captionof{figure}{A patient forum about myocardial infarction; French original, English translation.}
\label{fig:bt-web}
\end{center}
\medskip

Biomedical text is not a single language. The same medical knowledge changes with
who writes it, who reads it, and what the text is meant to do. Clinical notes,
research articles, and patient-facing texts therefore offer different views of
medicine, and only the latter two are widely available. The challenge for the rest
of this thesis is to use this public text without losing sight of the clinical
language we ultimately need.
